\tocsection{\textbf{9/15/26 Lecture 6:} Exponential Families}

\subsection*{Exponential Families}
\sectionspace

\begin{definition}[\textbf{Definition 6.1:} Exponential Families]
    An \emph{$s$-parameter exponential family} is a family $\cal{P} = \{P_{\eta} \col \eta \in \Xi\}$ with densities of the form
    \[ p_{\eta}(x) = e^{\eta'T(x) - A(\eta)}h(x) \]
    with respect to a common dominating measure $\mu$ on $X$. Recall that $\eta' := \eta^{\top}$. 
\end{definition}
Each of the components in the density has a name:
\begin{itemize}
    \item $\eta \in \Xi \subseteq \R^s$ is called the \emph{natural parameter}
    \item $T : \chi \to \R^{s}$ is the $s$-dimensional \emph{sufficient statistic}
    \item $h : \chi \to [0, \infty)$ is called the \emph{base-density} or \emph{carrier-density}
    \item $A : \Xi \to \R$ is called the \emph{log-partition function} or \emph{normalizing constant}
\end{itemize}
Now that, as $p_{\eta}$ must be a probability density, we can solve for $A(\eta)$ to get
\[ A(\eta) = \log \int_{\chi} e^{\eta'T(x)}h(x)d\mu \leq \infty \]
Notice that sometimes this can be infinity, which is not good, so we give the following definition.

\begin{definition}[\textbf{Definition 6.2:} Natural Parameter Space]
    The \emph{natural parameter space} is $\Xi_{1} = \{\eta \col A(\eta) < \infty\}$. That is, the set of $\eta$ for which $A(\eta)$ is finite.
\end{definition}

It can be shown that (using H\"older's inequality!) as the function $A(\eta)$ is convex, we also know $\Xi_1$ is a convex set, which may come in handy.\\

\textbf{Heuristically}: The reason we care about these families is that they are in some sense the closest we can get to a ``hyperplane'' of densities. The reason is that, absorbing $h(x)$ into $\mu$, we get $p_{\eta}(x) = e^{\eta'T(x) - A(\eta)}$, so
\[ \log p_{\eta}(x) = \eta'T(x) - A(\eta) \]
removing the normalizing factor, we get a hyperplane! Now, there is still the log in front, but this turns out to be more of a blessing in the curse, because logs turn multiplication into addition, so we are, in a sense, linearizing multplication which is good since we do a lot more multiplication than addition in statistics. 

\begin{example}{\textbf{Exapmle 1}: Poisson}
    Let $X \sim \Pois(\lambda)$ which has pmf $\frac{\lambda^{x}e^{-\lambda}}{x!}$. Then, we can write
    \[ p_{\lambda}(x) = e^{x\log\lambda - \lambda} \cdot \frac{1}{x!} \]
    Now, to put it in the form of the exponential family above, we must then have that $\eta(\lambda) = \log(\lambda)$, $T(x) = x$, $A(\eta) = \lambda = e^{\eta}$, and $h(x) = \frac{1}{x!}$, so that
    \[ p_{\eta}(x) = e^{\eta x - e^{\eta}} \frac{1}{x!} \]
    A question one may have is that: Why didn't we just leave it in terms of $\lambda$? Well, having it explicitly in this form will prove itself useful soon.
\end{example}

\begin{example}{\textbf{Example 2:} Poisson, n Variables}
    Let $X_1, \ldots, X_n \iidsim \Pois(\lambda)$. We see that
    \begin{align*}
        p_{\lambda}(X) &= \prod_{i = 1}^{n} e^{\log \lambda x_i - \lambda} \frac{1}{x_i!}\\
        &= e^{\log \lambda \sum_{i = 1}^{n}x_i - n\lambda} \cdot \frac{1}{\prod_{i = 1}^{n} x_i!}
    \end{align*}
    Thus, in this case $\eta = \log \lambda$, $T(x) = \sum_{i = 1}^{n} x_i$, $A(\eta) = ne^{\eta}$, and $h(x) = \frac{1}{\prod_{i = 1}^{n} x_i!}$
\end{example}

\begin{example}{\textbf{Example 3:} Generic, n Variables}
    Suppose that 
    \[ X_1, \ldots, X_n \iidsim p_{eta}^{(1)} = e^{\eta'T^{(1)}(x) - A^{(1)}(\eta)} \cdot h^{(1)}(x) \]
    where the $(1)$ means each density is for a single variable. Then,
    \begin{align*}
        p_{\eta}(x) &= \prod_{i = 1}^{n} e^{\eta'T^{(1)}(x_i) - A^{(1)}(\eta)}h^{(1)}(x)\\
        &= e^{\eta'\sum_{i = 1}^{n} T^{(1)}(x_i) - nA^{(1)}(\eta)}\prod_{i =1}^{n} h^{(1)}(x_i)
    \end{align*}
    Thus, we see that $\eta = \eta$, $T(x) = \sum_{i = 1}^{n}T^{(1)}(x_i)$, $A(\eta) = nA^{(1)}(\eta)$, and $h(x) = \prod_{i = 1}^{n} h^{(1)}(x_i)$. 
\end{example}

\subsection*{Differential Identites}
\sectionspace

Notice that on $\Xi_1$
\[ e^{A(\eta)} = \int e^{\eta'T(x)}h(x)d\mu \]
Let $g(\eta) = \int f(x)e^{\eta'T(x)}h(x)d\mu$. On the interior of where $\int \abs{f(x)e^{\eta'T(x)}h(x)}d\mu$ is finite, $g(\eta)$ is infinitely differentiable (this is via a theorem from real analysis), so on the interior of $\Xi_1$, we see that $e^{A(\eta)}$ is infinitely differentiable. Thus, taking the derivative with respect to a component $\eta_j$, we see that, assuming we can differentiate under the integral sign. 
\begin{align*}
    \frac{\del}{\del \eta_j} e^{A(\eta)} &= \int \frac{\del}{\del \eta_j} e^{\eta'T(x)}h(x)d\mu\\
    \frac{\del A}{\del \eta_j}(\eta)e^{A(\eta)} &= \int T_{j}(x)e^{\eta'T(x)}h(x)d\mu\\
    \frac{\del A}{\del \eta_j}(\eta) &= \int T_{j}(x)e^{\eta'T(x) - A(\eta)}h(x)d\mu\\
    &= \bb{E}_{\eta}[T_{j}(X)]
\end{align*}
Thus, as this holds for all partial derivatives, we see that $\nabla_{\eta}A(\eta) = \bb{E}_{\eta}\left[T(X)\right]$. That's a nice way to compute the expectation, but can we get more?\\

Taking the second derivative with respect to another component $\eta_k$ (which may be the same as before), we see
\begin{align*}
    \frac{\del^2}{\del \eta_j \del \eta_k}e^{A(\eta)} &= \int \frac{\del^2}{\del \eta_j \del \eta_k} e^{\eta'T(x)}h(x)d\mu\\
    e^{A(\mu)}\left(\frac{\del^2A}{\del \eta_j\del\eta_k}(\eta) + \frac{\del A}{\del \eta_j}(\eta)\frac{\del A}{\del \eta_j}(\eta)\right) &= \int T_{j}(x)T_{k}(x)e^{\eta'T(x)}h(x)d\mu\\
    \left(\frac{\del^2A}{\del \eta_j\del\eta_k}(\eta) + \frac{\del A}{\del \eta_j}(\eta)\frac{\del A}{\del \eta_j}(\eta)\right) &= \int T_{j}(x)t_{k}(x)e^{\eta'T(x) - A(\eta)}h(x)d\mu\\
    \left(\frac{\del^2 A}{\del\eta_j\del\eta_k}(\eta) + \bb{E}_{\eta}[T_{j}(X)]\bb{E}_{\eta}[T_{k}(X)]\right) &= \mathbb{E}_{\eta}[T_j(X)T_k(X)]\\
    \frac{\del^2A}{\del \eta_j \del \eta_k}(\eta) &= \tt{Cov}_{\eta}(T_{j}(X), T_{k}(X))
\end{align*}
Thus, we see that $\nabla_{\eta}^2 A(\eta) = \tt{Var}_{\eta}(T(X)) \in \R^{s \times s}$, where $\nabla_{\eta}^2$ is the Hessian, and $\tt{Var}_{\eta}(T(X))$ is the variance-covariance matrix.\\

\begin{example}{\textbf{Example 4:} Poisson}
    Recall from above we have that $T(x) = x$ and $A(\eta) = e^{\eta} = \lambda$. Then, we see that
    \[ \bb{E}_{\eta}[X] = \frac{dA}{d\eta} = e^{\eta} = \lambda \]
    and
    \[ \tt{Var}_{\eta}[X] = \frac{d^{2}A}{d\eta^2} = e^{\eta} = \lambda \]
    which is expected! Notice that if we were to not replace $\eta$ with $\lambda$ as we did above, and had $A(\lambda) = \lambda$, then we would see that
    \[ \bb{E}_{\eta}[X] = 1 \quand \tt{Var}_{\eta}[X] = 0 \]
    which is wrong.
\end{example}
\clearpage

\subsection*{Moment Generating Functions}

Now, because taking derivatives gives moments, there must be some connection to moments right? Here it is:

\begin{align*}
    M_{\eta}^{T(X)}(u) &= \bb{E}_{\eta}[e^{u'T(x)}]\\
    &= \int e^{u'T(x)}e^{\eta'T(x) - A(\eta)}h(x)d\mu\\
    &= \frac{e^{A(\eta + u)}}{e^{A(\eta)}}\int e^{(u + \eta)'T(x) - A(u + \eta)}h(x)d\mu(x)\\
    &= e^{A(\eta + u) - A(\eta)}
\end{align*}
where the third step follows from the fact we are integrating a probability density over the space, which requires that $u + \eta \in \xi_1$, but forget about it. \\

Anyways, taking derivatives with respect to $u$ gives
\[ \nabla_{u} e^{A(\eta + u) - A(\eta)} = \nabla(A(\eta + u))e^{A(\eta + u) - A(\eta)} \]
Then, as $u = 0$, we get
\[ \nabla_{u} M_{\eta}^{T(X)}(u) = \nabla(A(\eta)) = \E_{\eta}[T(X)] \]
which is as expected!\\

There is also another related function called the cumulant-generating function, which looks like
\[ K_{\eta}^{T(X)}(u) = \log M_{\eta}^{T(X)}(u) = A(\eta + u) - A(\eta) \]
but we will talk about these later. 
